\section{比较方法在本模型下的准则}\label{app:baselines}
\subsection{单源谱、MUSIC及MVDR-known}
LS取峰使用式\eqref{eq:single}。对MUSIC和MVDR，将同一记录划分为$S$个不重叠片段，各片段长度为$N_s=N/S$。在片段DFT频率轴上选择相同物理频带，构造
\begin{equation}
 \widehat{\bm C}_f=S^{-1}\sum_{b=1}^S\bm z_{bf}\bm z_{bf}^H .
\end{equation}
统计样本均来自该条观测，不从其他试次借快拍。完整记录与片段DFT分别保持$1/N$和$1/N_s$归一化；其白噪声方差分别为$\sigma_t^2/N$和$\sigma_t^2/N_s$。

设$\bm E_{n,f}$为$\widehat{\bm C}_f$最小$L-K$个特征值对应的特征向量。本文MUSIC适配版使用
\begin{align}
 S_{\mathrm{MU}}(\bm p)&=\left[\sum_f w_f
 \bm q_f(\bm p)^H\bm E_{n,f}\bm E_{n,f}^H\bm q_f(\bm p)\right]^{-1},\label{eq:music}\\
 \widetilde w_f&=\max\{\tr(\widehat{\bm C}_f)-L\sigma_t^2/N_s,0\},\quad
 w_f=\widetilde w_f/\sum_v\widetilde w_v .\nonumber
\end{align}
所有权重为零时退回均匀权重。已知$K$用于划分子空间，$w_f$由观测计算。处理频带外的频点不参与子空间分解。

MVDR-known采用
\begin{align}
 \bm C_{f,\delta}&=\widehat{\bm C}_f+
 \delta\,\tr(\widehat{\bm C}_f)\bm I/L,\nonumber\\
 S_{\mathrm{MV}}(\bm p)&=F_s^{-1}\sum_f
 [\bm q_f(\bm p)^H\bm C_{f,\delta}^{-1}\bm q_f(\bm p)]^{-1} .\label{eq:mvdr}
\end{align}
这里$\delta$为相对加载量，$F_s$为片段频率轴上的有效频点数。“known”表示采用本文已知相对传播响应；命名不等同于原DPD-HR公式的逐项复现。

\subsection{充分AP、多初值联合ML与GA}
充分AP对每个源反复执行全域条件扫描、候选比较及连续精修，更新时采用式\eqref{eq:profile}。它与短流程MCJ-AP共享物理目标，区别在于全域核查和初值处理的次数。完整AP保留早期有效解，因而比较成本时可区分首次可靠输出与最终核查结束。

多初值联合ML从四个唯一的谱峰组合和四个Sobol覆盖组合出发，以有界L-BFGS-B最小化$J$，输出最小残差解。候选组合的无序重复在初始化阶段删除。GA在相同四维位置目标上进行锦标赛选择、模拟二进制交叉和多项式变异\cite{DebAgrawal1995}，保留4个精英，最后作联合连续精修。锦标赛规模为3，交叉概率为0.9、分布指数为15，每维变异概率为0.25、分布指数为20。GA的总体成本由种群、代数和实际目标调用共同决定。

\subsection{多频共享功率SBL}
在公共位置字典$\bm D_f=[\bm q_f(\bm g_1),\ldots,\bm q_f(\bm g_G)]$上，令逐频系数先验为$\bm z_f\sim\mathcal{CN}(0,\bm\Gamma)$，$\bm\Gamma=\diag(\bm\gamma)$跨频共享\cite{Nannuru2019}。于是
\begin{align}
 \bm C_f&=\sigma_f^2\bm I+\bm D_f\bm\Gamma\bm D_f^H,\label{eq:sblCov}\\
 \bm\mu_f&=\bm\Gamma\bm D_f^H\bm C_f^{-1}\bm x_f,\nonumber\\
 \mathcal E(\bm\gamma)&=F^{-1}\sum_f[\log\det\bm C_f+\bm x_f^H\bm C_f^{-1}\bm x_f] .\label{eq:sblEvidence}
\end{align}
平坦频谱先验下，半步固定点更新写为
\begin{equation}
 \gamma_j^+=\gamma_j\left[
 \frac{\sum_f|\bm d_{fj}^H\bm C_f^{-1}\bm x_f|^2}
 {\sum_f\bm d_{fj}^H\bm C_f^{-1}\bm d_{fj}}
 \right]^{1/2} .\label{eq:sblUpdate}
\end{equation}
证据回溯控制更新，每5轮对8个活动字典点作离网优化。每维位移限制为20\,m，坐标尺度为1000\,m，L-BFGS-B最多20轮，最终从单点$\gamma_j$取峰。每5轮检查一次：证据归一化变化小于$10^{-5}$，$\gamma$的L1相对变化小于1\%，输出位置和活动点最大移动均小于0.05\,m；连续两次检查满足时停止，否则最多300轮。谱峰分组候选不并入本文提出方法。

\subsection{局部精修与停止参数}
AP的交替和联合版本在归一化目标$J$上使用相同投影梯度阈值$10^{-8}$，目标差停止项设为零。每次L-BFGS-B最多100轮，线搜索最多25次；交替版本最多12个完整位置扫描。初始化的单源精修采用目标相对容差$10^{-12}$。候选在$3\times3$邻域中取局部极大值，再按30\,m间隔保留；该间隔是候选去重参数，不是成功门槛。

VP-BNOMP每次局部Newton最多5步，Hessian最小特征值下限为$10^{-9}$，位置步长上限为50\,m。按$1,1/2,1/4,1/8,1/16$依次回溯，接受下降步，局部小步阈值为0.01\,m。循环中目标下降小于$10^{-8}$且最大位置移动小于0.05\,m，连续两轮满足时停止。已完成配对组的循环上限为12；扩大比较固定为200，以分别记录原生停止与预算触顶。

\subsection{固定系数BNOMP对应版与消元对照}
原文对应版根据给定系数$q_f^{(0)}$使用
\begin{equation}
 L_{\mathrm{FC}}(\bm p)=E_x^{-1}\sum_f
 \|\bm y_f-\bm a_f(\bm p)q_f^{(0)}\|^2 .\label{eq:fixedcoef}
\end{equation}
对位置求Newton方向，再按该版的时序重估系数。本文以二维多频块NOMP\cite{Yang2022}为对应来源，移植的是系数固定下的位置更新及循环结构，观测由声学波束形成改为本研究的四站时延--幅度响应。符号、单位、边界和数值保护均按本模型登记。

受控消元版只把式\eqref{eq:fixedcoef}改成式\eqref{eq:localprofile}，用于隔离曲率效应；VP-BNOMP还采用其自身的共同系数重估时序。受控版的倍率与完整算法的对比结果分开报告，避免把多个时序差异都归于一处。

\section{消元局部导数与实现核验}\label{app:derivatives}
令$\kappa=2\pi f/c_0$，$\ell_l=\log a_l=-\log d_l-\mathrm j\kappa d_l$。在远离站点的区域，
\begin{align}
 \nabla\ell_l&=(-d_l^{-1}-\mathrm j\kappa)\bm v_l,\\
 \nabla^2\ell_l&=d_l^{-2}\bm v_l\bm v_l^T
 +\frac{-d_l^{-1}-\mathrm j\kappa}{d_l}
 (\bm I-\bm v_l\bm v_l^T),\\
 \nabla^2 a_l&=a_l[\nabla^2\ell_l+(\nabla\ell_l)(\nabla\ell_l)^T].\label{eq:responseHessian}
\end{align}
对固定的$\bm y_f$，将$q_f$拆为实部和虚部，其系数Hessian为
\begin{equation}
 H_{uu,f}=2\|\bm a_f\|^2\bm I_2/E_x .
\end{equation}
逐频计算混合导数并累加Schur补，即可得到式\eqref{eq:localprofile}的$2\times2$实Hessian。实现核验分别对照直接最小二乘、投影消元、解析梯度、有限差分和小规模直接矩阵。理论恒等式核验与定位性能实验分别记录。
